% =========================================================================
% ANEXOS DEL PLAN DE TESIS
% =========================================================================
\clearpage
\phantomsection
\chapter*{Anexos}
\addcontentsline{toc}{chapter}{Anexos}

\section*{Anexo 1: Matriz de consistencia lógica}
\phantomsection
\label{anexo:matriz_consistencia}
\addcontentsline{toc}{section}{Anexo 1: Matriz de consistencia lógica}

\begin{table}[H]
\centering
\begin{threeparttable}
\caption{Matriz de consistencia lógica de la investigación}
\label{tab:anexo_matriz_consistencia}
\scriptsize
\begin{tabularx}{\textwidth}{@{}p{2.6cm}p{2.6cm}p{2.6cm}p{2.8cm}X@{}}
\toprule
\rowcolor{tcAzul!12}
\textbf{Problemas} & \textbf{Objetivos} & \textbf{Hipótesis} & \textbf{Variables e indicadores} & \textbf{Método} \\
\midrule
\textbf{Problema general:} ¿Cómo se comporta una arquitectura SLM--RAG para la generación contextualizada en Quechua Collao dentro de un entorno computacional sin conexión, y qué aporte adicional presenta QLoRA si el ajuste y la inferencia con adaptador son viables? &
\textbf{Objetivo general:} Diseñar e implementar una arquitectura SLM--RAG para la generación contextualizada en Quechua Collao y evaluar el aporte adicional de QLoRA si el ajuste y la inferencia con adaptador son viables. &
\textbf{Hipótesis general:} Si QLoRA es viable, SLM--RAG--QLoRA obtendrá mayor fidelidad al contexto y similitud con referencias que SLM--RAG; también se describirán sus diferencias frente al SLM base y su costo computacional. &
\textbf{VI:} configuración del sistema. \textbf{VD:} similitud con referencias, fidelidad y eficiencia. &
Enfoque cuantitativo; investigación aplicada y tecnológica; experimento computacional comparativo con consultas repetidas. \\
\addlinespace
\textbf{Problema específico 1:} ¿Cómo estructurar e indexar una base documental bilingüe para recuperar información pertinente en Quechua Collao? &
\textbf{Objetivo específico 1:} Curar, estructurar e indexar documentos bilingües para recuperación local. &
\textbf{Hipótesis específica 1:} La indexación semántica recuperará una mayor proporción de fragmentos pertinentes que la búsqueda léxica de referencia, según las consultas y etiquetas definidas. &
Precisión de recuperación (P@k), cobertura de fuentes y trazabilidad documental. &
Corpus con etiquetas de idioma y variedad; índice local y consultas con evidencia relevante identificada. \\
\addlinespace
\textbf{Problema específico 2:} Si el ajuste y la inferencia con adaptador son viables, ¿qué diferencia produce QLoRA en la fidelidad al contexto de SLM--RAG bajo las mismas condiciones locales? &
\textbf{Objetivo específico 2:} Integrar QLoRA en SLM--RAG y comparar su fidelidad al contexto con SLM--RAG sin adaptación, si el ajuste y la inferencia con adaptador son viables. &
\textbf{Hipótesis específica 2:} Si el ajuste QLoRA y la inferencia con el adaptador son viables, SLM--RAG--QLoRA tendrá mayor fidelidad al contexto que SLM--RAG, manteniendo constantes las demás condiciones. &
Fidelidad al contexto; latencia; tokens por segundo; RAM y VRAM. &
Registro de configuración y telemetría; prueba de viabilidad previa al experimento completo. \\
\addlinespace
\textbf{Problema específico 3:} ¿Qué diferencias de similitud de respuesta, fidelidad al contexto y eficiencia se observan entre SLM, SLM--RAG y SLM--RAG--QLoRA, cuando esta última condición sea viable? &
\textbf{Objetivo específico 3:} Comparar SLM base, SLM--RAG y SLM--RAG--QLoRA, cuando esta última condición sea viable, mediante indicadores de similitud de respuesta, fidelidad al contexto y eficiencia computacional. &
\textbf{Hipótesis específica 3:} Si las tres configuraciones pueden ejecutarse, SLM--RAG--QLoRA obtendrá mayor similitud con las respuestas de referencia y mayor fidelidad que SLM y SLM--RAG. &
chrF++; BLEU; P@k; fidelidad al contexto; latencia y memoria. &
Banco provisional de consultas; comparaciones pareadas y prueba global si se evalúan tres condiciones. \\
\bottomrule
\end{tabularx}
\begin{tablenotes}
\scriptsize
\item \textit{Nota.} El tamaño de 100 consultas y la evaluación de QLoRA son provisionales. Se confirmarán después de validar el corpus y la viabilidad del hardware.
\end{tablenotes}
\end{threeparttable}
\end{table}

\clearpage
\section*{Anexo 2: Ficha técnica de registro experimental}
\phantomsection
\label{anexo:ficha_tecnica}
\addcontentsline{toc}{section}{Anexo 2: Ficha técnica de registro experimental}

\begin{table}[H]
\centering
\caption{Ficha para registrar una corrida experimental}
\label{tab:ficha_tecnica}
\begin{tabularx}{\textwidth}{@{}p{5.0cm}X@{}}
\toprule
\textbf{Campo} & \textbf{Registro durante la ejecución} \\
\midrule
Identificador & Completar al ejecutar \\
Configuración & SLM base / SLM--RAG / SLM--RAG--QLoRA, según viabilidad \\
Dispositivo y entorno & Modelo del equipo, sistema operativo y versiones \\
Modelo y cuantización & Nombre, versión y formato realmente utilizados \\
Consulta y variedad & Consulta evaluada, idioma y variedad identificados \\
Fragmentos recuperados & Identificadores de fuente y fragmentos entregados al modelo \\
Latencia y rendimiento & Medición en milisegundos y tokens por segundo \\
Consumo de memoria & Pico de RAM y VRAM medidos \\
Salida y métricas & Respuesta generada y puntuaciones calculadas \\
Observaciones & Errores, interrupciones y notas del evaluador \\
\bottomrule
\end{tabularx}
\end{table}

\clearpage
\section*{Anexo 3: Rúbrica de revisión lingüística y factual}
\phantomsection
\label{anexo:revision_linguistica}
\addcontentsline{toc}{section}{Anexo 3: Rúbrica de revisión lingüística y factual}

\begin{table}[H]
\centering
\begin{threeparttable}
\caption{Rúbrica para revisar respuestas generadas en Quechua Collao}
\label{tab:revision_linguistica}
\small
\begin{tabularx}{\textwidth}{@{}p{2.7cm}X X X@{}}
\toprule
\textbf{Criterio} & \textbf{Insuficiente (1)} & \textbf{Parcial (2--3)} & \textbf{Adecuado (4--5)} \\
\midrule
Comprensión de la consulta & No responde a la pregunta o altera su sentido. & Responde parcialmente o añade contenido no solicitado. & Responde al propósito de la consulta de forma directa. \\
\addlinespace
Corrección lingüística & Errores que impiden comprender el sentido en la variedad indicada. & El sentido general se entiende, aunque requiere correcciones. & Redacción comprensible y consistente con Quechua Collao. \\
\addlinespace
Fidelidad a la fuente & Incluye afirmaciones que contradicen o no aparecen en la evidencia. & La mayoría de afirmaciones se sustenta, con omisiones o dudas. & Las afirmaciones evaluables se sustentan en la fuente proporcionada. \\
\addlinespace
Variedad y vocabulario & Mezcla variedades o usa vocabulario que cambia el significado. & Mantiene parcialmente la variedad, con usos que requieren revisión. & Conserva la variedad y el vocabulario pertinentes a la consulta. \\
\bottomrule
\end{tabularx}
\begin{tablenotes}
\scriptsize
\item \textit{Nota.} La ficha será aplicada por una persona competente en Quechua Collao, con criterios acordados antes de evaluar. Si participan evaluadores, se registrará el procedimiento de revisión y resolución de discrepancias. La rúbrica complementa, pero no sustituye, las métricas automáticas.
\end{tablenotes}
\end{threeparttable}
\end{table}

\clearpage
\section*{Anexo 4: Plantilla del banco de consultas de evaluación}
\phantomsection
\label{anexo:banco_consultas}
\addcontentsline{toc}{section}{Anexo 4: Plantilla del banco de consultas de evaluación}

\begin{table}[H]
\centering
\begin{threeparttable}
\caption{Estructura prevista para documentar consultas de evaluación}
\label{tab:banco_consultas}
\small
\begin{tabularx}{\textwidth}{@{}c p{3.1cm} X p{4.1cm}@{}}
\toprule
\textbf{ID} & \textbf{Tema e idioma / variedad} & \textbf{Consulta} & \textbf{Fuente, referencia y pasajes pertinentes} \\
\midrule
Q01 & Pendiente de selección; por identificar & Completar tras curar las fuentes & Registrar documento, respuesta verificable e identificadores de evidencia \\
\addlinespace
Q02 & Pendiente de selección; por identificar & Completar tras curar las fuentes & Registrar documento, respuesta verificable e identificadores de evidencia \\
\addlinespace
Q03 & Pendiente de selección; por identificar & Completar tras curar las fuentes & Registrar documento, respuesta verificable e identificadores de evidencia \\
\bottomrule
\end{tabularx}
\begin{tablenotes}
\scriptsize
\item \textit{Nota.} Plantilla del instrumento; no representa consultas ni resultados ya validados. Se registrarán los fragmentos considerados pertinentes para contrastar la recuperación. El banco previsto se completará después de verificar las fuentes y la cobertura lingüística.
\end{tablenotes}
\end{threeparttable}
\end{table}
